\documentclass[aspectratio=169,10pt]{beamer}

\usepackage[utf8]{inputenc}
\usepackage[T1]{fontenc}
\usepackage[spanish,es-noshorthands,es-tabla]{babel}
\usepackage{booktabs}
\usepackage{tikz}
\usetikzlibrary{positioning,arrows.meta}
\usepackage{listings}
\usepackage{xcolor}

\usetheme{Madrid}
\setbeamertemplate{navigation symbols}{}

\lstset{
  basicstyle=\ttfamily\scriptsize,
  keywordstyle=\bfseries,
  showstringspaces=false,
  frame=single,
  rulecolor=\color{gray!40},
  literate={á}{{\'a}}1 {é}{{\'e}}1 {í}{{\'i}}1 {ó}{{\'o}}1 {ú}{{\'u}}1 {ñ}{{\~n}}1
}

\newcommand{\hecho}{\textbf{Hecho}}
\newcommand{\pendiente}{Pendiente}

\title[Civic Security]{Civic Security}
\subtitle{Avance 1 del MVP: arquitectura y base de datos}
% Completar con los nombres de los tres integrantes del equipo
\author[Equipo Civic Security]{Equipo Civic Security}
\institute[UCSP]{ET201 -- Formación de Empresas de Base Tecnológica I \\ Universidad Católica San Pablo}
\date[28/10/2026]{28 de octubre de 2026}

\begin{document}

\begin{frame}
  \titlepage
\end{frame}

% ------------------------------------------------------------------
\begin{frame}{El problema y nuestra propuesta}
  \begin{columns}[T]
    \begin{column}{0.48\textwidth}
      \textbf{Problema validado en entrevistas}
      \begin{itemize}
        \item La información de seguridad circula en grupos de WhatsApp, pero \alert{nadie confía en ella} porque no está verificada.
        \item Muchas personas \alert{no reportan por miedo} a ser identificadas.
      \end{itemize}
    \end{column}
    \begin{column}{0.48\textwidth}
      \textbf{Civic Security}
      \begin{enumerate}
        \item Reporte \textbf{anónimo} y seguro.
        \item \textbf{Validación comunitaria}: 3 confirmaciones de cuentas verificadas a menos de 1\,km.
        \item \textbf{Mapa de calor} y alertas por zona en tiempo real.
      \end{enumerate}
    \end{column}
  \end{columns}
\end{frame}

% ------------------------------------------------------------------
\begin{frame}{Qué logramos en este avance}
  \begin{enumerate}
    \item \textbf{Arquitectura definida y documentada} (v0.2), con las decisiones del equipo incorporadas.
    \item \textbf{Repositorio público} con historial de commits del equipo.
    \item \textbf{Base de datos funcionando}: PostgreSQL + PostGIS en Docker.
    \item \textbf{Backend Phoenix} con verificación de salud (\texttt{/health}).
    \item \textbf{Modelo de incidente} con las reglas de negocio y \textbf{10 pruebas automáticas} en verde.
  \end{enumerate}
  \vfill
  \small Repositorio: \url{https://github.com/aaabriceno/Proyecto-MVP-Civic-Security}
\end{frame}

% ------------------------------------------------------------------
\begin{frame}{Arquitectura: monolito modular}
  \centering
  \begin{tikzpicture}[
      caja/.style={draw=structure, fill=structure!10, rounded corners, minimum width=3.3cm, minimum height=1.2cm, align=center, font=\small},
      ext/.style={draw=gray, dashed, rounded corners, minimum width=2.6cm, minimum height=0.8cm, align=center, font=\scriptsize},
      flecha/.style={-{Stealth}, thick, structure}
    ]
    \node[caja] (app) {\textbf{App móvil}\\Flutter (Android)};
    \node[caja, right=1.6cm of app] (api) {\textbf{Backend}\\Elixir + Phoenix};
    \node[caja, right=1.6cm of api] (db) {\textbf{Base de datos}\\PostgreSQL + PostGIS};
    \node[ext, below=0.9cm of api] (fcm) {Notificaciones push\\(Firebase)};
    \node[ext, below=0.9cm of db] (r2) {Fotos\\(Cloudflare R2)};

    \draw[flecha] (app) -- node[above, font=\scriptsize]{REST + WebSocket} (api);
    \draw[flecha] (api) -- node[above, font=\scriptsize]{SQL} (db);
    \draw[flecha, gray] (api) -- (fcm);
    \draw[flecha, gray] (api) -- (r2);
  \end{tikzpicture}

  \vspace{0.6em}
  \begin{columns}[T]
    \begin{column}{0.48\textwidth}
      \small\textbf{Grano fino (mantenibilidad):} el código se divide en módulos por dominio con fronteras claras.
    \end{column}
    \begin{column}{0.48\textwidth}
      \small\textbf{Grano grueso (rendimiento):} se despliega como un solo servicio y una sola base de datos, sin la complejidad de microservicios.
    \end{column}
  \end{columns}
\end{frame}

% ------------------------------------------------------------------
\begin{frame}{Decisiones del equipo}
  \begin{table}
    \small
    \begin{tabular}{p{0.3\textwidth}p{0.62\textwidth}}
      \toprule
      \textbf{Tema} & \textbf{Decisión} \\
      \midrule
      Verificación de persona & DNI, guardando solo un \textbf{HMAC} (nunca el número). Una cuenta por DNI. \\
      Zonas de alerta & En el servidor: celda aproximada en claro y \textbf{centro exacto cifrado}. \\
      Red P2P & Descartada: sin servidor central no se puede garantizar la validación comunitaria. \\
      Plataforma de la demo & Solo Android. \\
      Fuera del MVP & Entrega cifrada a la PNP, eventos de tráfico, passkeys. \\
      \bottomrule
    \end{tabular}
  \end{table}
\end{frame}

% ------------------------------------------------------------------
\begin{frame}{Privacidad por diseño}
  \textbf{El anonimato se garantiza por la estructura de los datos, no solo por una política.}
  \vspace{0.5em}
  \begin{table}
    \small
    \begin{tabular}{lll}
      \toprule
      \textbf{Dato} & \textbf{Tratamiento} & \textbf{Por qué} \\
      \midrule
      Reporte & Sin campo de usuario ni IP & Nadie puede saber quién reportó \\
      DNI & HMAC con clave secreta & Irreversible; solo detecta duplicados \\
      Centro de zona & Cifrado AES-256-GCM & Se descifra solo en memoria \\
      Fotos & Rostros difuminados, sin EXIF & No exponen a terceros ni la ubicación \\
      \bottomrule
    \end{tabular}
  \end{table}
  \vfill
  \small Estándares de referencia: OWASP MASVS, OWASP ASVS, NIST SP 800-63B y Ley N.º 29733.
\end{frame}

% ------------------------------------------------------------------
\begin{frame}[fragile]{Base de datos: tabla \texttt{incidentes}}
  \begin{columns}[T]
    \begin{column}{0.55\textwidth}
\begin{lstlisting}
 Column          | Type
-----------------+-----------------------
 id              | uuid
 tipo            | character varying(255)
 descripcion     | character varying(500)
 ubicacion       | geography(Point,4326)
 estado          | character varying(255)
 idempotency_key | uuid
 inserted_at     | timestamp
 updated_at      | timestamp
Indexes:
 incidentes_pkey   PRIMARY KEY (id)
 ..._idempotency_key_index  UNIQUE
 ..._inserted_at_index      btree
 ..._ubicacion_index        gist
\end{lstlisting}
    \end{column}
    \begin{column}{0.42\textwidth}
      \small
      \begin{itemize}
        \item \textbf{UUID}: los reportes no se pueden enumerar (1, 2, 3\ldots).
        \item \textbf{geography}: distancias reales en metros, para la regla de 1\,km.
        \item \textbf{Índice GiST}: búsquedas rápidas por zona del mapa.
        \item \textbf{UNIQUE}: impide reportes duplicados.
        \item \alert{Ninguna columna identifica al usuario.}
      \end{itemize}
    \end{column}
  \end{columns}
\end{frame}

% ------------------------------------------------------------------
\begin{frame}[fragile]{Reglas de negocio en el código}
  \begin{columns}[T]
    \begin{column}{0.58\textwidth}
\begin{lstlisting}[language=Ruby,basicstyle=\ttfamily\fontsize{6.5}{8}\selectfont]
def create_changeset(incident, attrs) do
  incident
  |> cast(attrs, [:tipo, :descripcion, :lat, :lng])
  |> validate_required([:tipo, :lat, :lng])
  |> update_change(:descripcion, &trim_or_nil/1)
  |> validate_length(:descripcion, max: 500)
  |> validate_arequipa()
  |> put_ubicacion()
  |> unique_constraint(:idempotency_key)
end
\end{lstlisting}
    \end{column}
    \begin{column}{0.39\textwidth}
      \small
      \begin{itemize}
        \item Solo 5 tipos de incidente válidos.
        \item Ubicación obligatoria y \textbf{dentro de Arequipa}.
        \item Descripción opcional, máximo 500 caracteres.
        \item Todo reporte nace como \textbf{``Sin verificar''}.
        \item \textbf{Lista blanca} de campos: nadie puede marcar su propio reporte como verificado.
        \item \textbf{Idempotencia}: si la app reintenta por mala señal, no se duplica.
      \end{itemize}
    \end{column}
  \end{columns}
\end{frame}

% ------------------------------------------------------------------
\begin{frame}[fragile]{Evidencia de funcionamiento}
  \begin{columns}[T]
    \begin{column}{0.5\textwidth}
      \textbf{\small Base de datos (Docker)}
\begin{lstlisting}
$ docker compose ps
civic_db | postgis/postgis:17-3.5
         | Up (healthy) | 127.0.0.1:5432

$ SELECT postgis_version();
3.5 USE_GEOS=1 USE_PROJ=1 USE_STATS=1
\end{lstlisting}
      \scriptsize La base de datos solo acepta conexiones locales (\texttt{127.0.0.1}).
    \end{column}
    \begin{column}{0.5\textwidth}
      \textbf{\small Pruebas automáticas}
\begin{lstlisting}
$ mix test --trace test/civic
* crea un incidente válido (sin_verificar)
* rechaza tipo desconocido
* rechaza ubicación fuera de Arequipa
* exige lat y lng
* descripción con solo espacios -> nil
* descripción de 501 caracteres inválida
* misma idempotency_key -> mismo incidente

10 tests, 0 failures
\end{lstlisting}
    \end{column}
  \end{columns}
\end{frame}

% ------------------------------------------------------------------
\begin{frame}{Estado del plan del primer entregable funcional}
  \small
  \begin{table}
    \begin{tabular}{clc}
      \toprule
      \textbf{\#} & \textbf{Tarea} & \textbf{Estado} \\
      \midrule
      1 & Base de datos PostgreSQL + PostGIS en Docker & \hecho \\
      2 & Backend Phoenix con verificación de salud & \hecho \\
      3 & Modelo de incidente con reglas de negocio & \hecho \\
      4 & API para crear reportes (\texttt{POST /v1/incidents}) & \pendiente \\
      5 & API del mapa (\texttt{GET /v1/map/snapshot}) & \pendiente \\
      6 & Canal en tiempo real & \pendiente \\
      7--10 & App Flutter: mapa, reporte y pines en tiempo real & \pendiente \\
      \bottomrule
    \end{tabular}
  \end{table}
  \vfill
  \textbf{Meta del entregable:} un reporte creado en un teléfono Android aparece en el mapa de otro teléfono en tiempo real.
\end{frame}

% ------------------------------------------------------------------


\end{document}
